%!TEX root = ../template.tex
%%%%%%%%%%%%%%%%%%%%%%%%%%%%%%%%%%%%%%%%%%%%%%%%%%%%%%%%%%%%%%%%%%%%
%% abstract-en.tex
%% NOVA thesis document file
%%
%% Abstract in English
%%%%%%%%%%%%%%%%%%%%%%%%%%%%%%%%%%%%%%%%%%%%%%%%%%%%%%%%%%%%%%%%%%%%

\typeout{NT FILE abstract-en.tex}%

The Industry 4.0 paradigm relies on the convergence of \glspl{CPS} and \glspl{CPPS} to achieve agile and resilient production. Within this landscape, \glspl{DT} provide virtual representations capable of mirroring physical assets through real-time bidirectional telemetry. However, the engineering lifecycle of industrial \glspl{DT} remains burdened by high deployment barriers. Traditional development relies on proprietary, vendor-locked simulation software, fragmented 3D \gls{CAD} files that lack kinematics and sensor semantics, and tedious manual matching of \gls{IO} addresses, making continuous reconfiguration prohibitively labor-intensive and error-prone.

To overcome these limitations, this dissertation presents an open, modular, model-driven framework for the automated procedural generation and live synchronization of industrial \glspl{DT}. Developed under the \gls{DSR} methodology, the architecture spans four decoupled operational tiers: physical survey and low-polygon asset digitization (Layer~0), no-code semantic asset authoring with 3D raycasting and \gls{AASX} packaging (Layer~1), containerized package storage and dynamic discovery via Eclipse BaSyx~v2 microservices (Layer~2), and procedural runtime scene assembly, contact physics simulation, dual-mode execution (\gls{VC} and live telemetry), and kinematic drift correction in Unity 3D (Layer~3).

The proposed framework was empirically validated on a physical modular \gls{FMS} bench built with \textit{fischertechnik} mechatronic hardware and driven by an Industrial Shields \gls{PLC} running FreeRTOS, as well as through parametric stress testing scaling from 5 to 1000 conveyor modules. Experimental results substantiate cold-boot instantiation of the pilot bench in $2.79 \pm 0.04\,$s while maintaining steady-state rendering rates of $489.91 \pm 49.21\,$\gls{FPS} and tail latencies below $9.94\,$ms. Scalability testing confirms near-linear scaling up to 250 modules and stable real-time execution up to 500 modules ($93\,$\gls{FPS}), establishing model-driven \gls{AAS} standards as an effective foundation for next-generation industrial digital twins.

\keywords{
  Digital Twin \and
  Cyber-Physical Systems \and
  Asset Administration Shell \and
  Procedural Simulation \and
  Virtual Commissioning \and
  Industry 4.0
}